\begin{description}
\item[(a)] \textbf{True}. The polar rod of the sundial faces the hemisphere
  that contains the south cardinal point as the fundamental pole, so the
  sundial was designed for a location in the southern hemisphere of the Earth.
\item[(b)] \textbf{False}. This curve represents the trajectory of the tip of
  the rod's shadow throughout the day of the summer solstice in the hemisphere
  in which the sundial is located.
\item[(c)] \textbf{True}. This straight line, which is perpendicular to the
  meridian line, represents the trajectory of the tip of the rod's shadow
  throughout the day of either southern autumnal or vernal equinox.
\item[(d)] \textbf{False}. This set of radial lines corresponds to the true
  local solar time. The 12h line coincides with north--south line, which must
  be the case for the true solar noon.
\item[(e)] \textbf{True}. The dashed lines and the solid lines form an angle
  corresponding to the longitude difference between the sundial and the center
  of the time zone. Therefore, the dashed lines correspond to the true solar
  time of the center of the time zone. An analemma shaped figure around the
  line that corresponds to a true local solar time indicates the mean solar
  time.
\end{description}
